\documentclass[12pt]{article}
\usepackage{amsmath}
\usepackage{amssymb}
\usepackage[french]{babel}
\usepackage[latin1]{inputenc}
\usepackage[T1]{fontenc}
\usepackage[all]{xy}
\usepackage{graphicx}
\usepackage{epsf}

\newcommand{\cours}{MAT501 -- Fondements et histoire des math�matiques}


\begin{document}

\large{

\begin{description}

\item[Axiome de Zorn] {\ } \\
Tout ensemble ordonn� inductif poss�de au moins un �l�ment maximal

\vspace{2cm}

\item[Axiome de Hausdorff-Kuratowski] {\ } \\
Toute cha�ne d'un ensemble ordonn� $E$ est incluse dans une cha�ne
maximale de $E$.

\vspace{2cm}

\item[Axiome du choix - version 1] {\ } \\
Soient $I$ un ensemble non vide et, pour chaque $i \in I$, un
ensemble non vide $X_i$. Alors, il existe une application $\phi :
I \to \bigcup\limits_{i \in I} X_i$ telle que, pour tout $i \in
I$, $\phi(i) \in X_i$.

\vspace{2cm}

\item[Axiome du choix - version 2] {\ } \\
Si $I$ et tous les $X_i$ sont non vides, alors le produit
$\prod\limits_{i \in I} X_i$ est non vide.

\vspace{2cm}

\item[Axiome du bon ordre] {\ } \\
Tout ensemble peut �tre muni d'un bon ordre.

\end{description}
}

\end{document}
